\section{Reachable sections and a cumulative volume bound}
\label{sec:volume}
We first treat a general finite collection of linear mean updates.
Let $\mathcal L$ be a finite set of linear maps of $\R^d$, containing
$I_d$. A finite seed alphabet has associated vectors $z(x)$, whose
convex hull $S$ has nonempty interior. Assume that every product of at
most $N$ allowed maps sends the seed vectors into $D=[-1,1]^d$.
After reading $x$ and $N$ commands $L_1,\ldots,L_N\in\mathcal L$,
a terminal query $j\in[d]$ requires a binary output of mean
$e_j^TL_N\cdots L_1z(x)$. Denote this complete array by $C_N$.

\begin{lemma}[The reachable-section lift]\label{lem:lift}
Every realization of $C_N$ with profile $(K_t)_{t=0}^N$ determines
full-dimensional polytopes $P_t\subset D$ such that
\begin{equation}\label{eq:lift-conditions}
 S\subset P_t,\qquad
 |\ext P_t|\le2^{K_t},\qquad
 LP_t\subset P_{t+1}\quad(L\in\mathcal L,\ t<N).
\end{equation}
The conclusion permits arbitrary hidden states and does not require
individual state continuations to belong to the residual affine span.
\end{lemma}
\begin{proof}
For each actual prefix $h$ at cut $t$, let $E_h\in\Delta_{K_t}$ be
its hidden state distribution. Put
\[
 H_t=\aff\{E_h:h\text{ a prefix at }t\},\qquad
 Q_t=H_t\cap\Delta_{K_t}.
\]
For any state distribution $e\in\Delta_{K_t}$, run the remaining
commands all equal to $I_d$. The vector of its $d$ terminal means is a
linear function $\pi_t(e)$ and belongs to $D$. Correctness on actual
prefixes gives $\pi_t(E_h)=z_h$, the specified predictive vector.
Set $P_t=\pi_t(Q_t)$.

Write $T_{t,L}$ for the realizing transition matrix of command $L$.
The identity $E_hT_{t,L}=E_{hL}$ implies
$H_tT_{t,L}\subset H_{t+1}$. Stochasticity also implies
$\Delta_{K_t}T_{t,L}\subset\Delta_{K_{t+1}}$; consequently
$Q_tT_{t,L}\subset Q_{t+1}$. On every actual prefix,
\[
 \pi_{t+1}(E_hT_{t,L})=L\pi_t(E_h).
\]
Both sides are affine in $E_h$, so this equality extends to all of $H_t$,
including its points with some negative coordinates. Restricting back to
$Q_t$ yields $LP_t\subset P_{t+1}$. Taking the idle command shows
$P_t\subset P_{t+1}$. The initial section contains all seeds, hence so
does every $P_t$.

It remains to bound vertices without assuming a minimal hidden
representation. In its affine hull, $Q_t$ is a bounded polytope defined
by the $K_t$ coordinate inequalities $e_s\ge0$. Two distinct vertices
cannot have the same zero-coordinate set. Indeed, if they did, their
difference would give a nonzero direction preserving all active
equalities, and sufficiently small motion in both directions from
either point would preserve the remaining strict inequalities. That
point would not be a vertex. There are at most $2^{K_t}$ zero sets.
A linear image of a polytope is the convex hull of the images of its
vertices, proving $|\ext P_t|\le|\ext Q_t|\le2^{K_t}$.
\end{proof}

The section $Q_t$, rather than the entire hidden simplex, is essential.
For an arbitrary unit hidden state, the identity relating the two probe
maps need not hold. It holds on the affine span of reachable
\emph{distributions}. This is the finite-horizon version of the
accessible-space/observable-quotient precaution in classical realization
\cite[Section 4]{MW}; the vertex bound accounts for the resulting lift.
In particular, replacing $2^{K_t}$ by $K_t$ without further hypotheses
would not be justified.

For a full-dimensional compact seed hull $S\subset D$, define
\begin{equation}\label{eq:expansion}
 \delta_M(S,D,\mathcal L)=
 \min_{\substack{S\subset P\subset D\\P\text{ a polytope},\ |\ext P|\le M}}
 \left\{\vol_d\left(\conv\bigcup_{L\in\mathcal L}LP\right)
                  -\vol_d(P)\right\}.
\end{equation}
If the feasible set is empty, put $\delta_M=+\infty$.
The maps in the union need not send all of $D$ into $D$.
Only the implementing sequence from Lemma~\ref{lem:lift} must obey that
constraint. The quantity is a geometric optimization on observable
polytopes, not an optimization over candidate transition tables.

\begin{theorem}[Volume budget for a state profile]\label{thm:volume}
Suppose $|\det L|=1$ for every $L\in\mathcal L$ and $I_d\in\mathcal L$.
If there is no full-dimensional polytope $P$ with
$S\subset P\subset D$ and $LP=P$ for all $L\in\mathcal L$, then every
nonempty minimum in \eqref{eq:expansion} is strictly positive.
Every realization of $C_N$ satisfies
\begin{equation}\label{eq:volume-budget}
 \sum_{t=0}^{N-1}\delta_{2^{K_t}}(S,D,\mathcal L)
                       \le\vol_d(D)-\vol_d(S).
\end{equation}
In particular, a uniform state bound $K$ implies
$N\delta_{2^K}\le\vol_d(D)-\vol_d(S)$.
\end{theorem}
\begin{proof}
The set of feasible polytopes with at most $M$ vertices is compact in
the Hausdorff metric: represent them by $M$ points in the compact set
$D$, permitting repetitions, and impose the closed condition that their
convex hull contain $S$. Convex hull, finite union followed by convex
hull, and volume of compact convex sets are continuous. Thus the
minimum is attained whenever the set is nonempty.

Because the identity is allowed, the displayed volume difference is
nonnegative. If it vanishes at $P$, the full-dimensional convex bodies
$P$ and $\conv\bigcup_LLP$ have equal volume and one contains the
other, so they are equal. Thus $LP\subset P$. The determinant hypothesis
makes their volumes equal, hence $LP=P$ for every $L$, contrary to the
assumption. Therefore $\delta_M>0$.

For the polytopes in Lemma~\ref{lem:lift},
\[
 \vol_d(P_{t+1})-\vol_d(P_t)
 \ge\vol_d\left(\conv\bigcup_LLP_t\right)-\vol_d(P_t)
 \ge\delta_{2^{K_t}}.
\]
Summing telescopes. The endpoint inclusions give
\eqref{eq:volume-budget}. The minima are nonincreasing in $M$, which
gives the assertion for a uniform bound.
\end{proof}

The volume hypothesis is a property of the mean dynamics and the
observation domain, not a small-ball or metric-dimension assumption.
Under an invertible affine change of predictive coordinates, using the
conjugate affine maps, both sides of \eqref{eq:volume-budget} are
multiplied by the same absolute determinant. Thus the normalized budget
is independent of this coordinate choice. The proof also works for
volume-preserving affine maps with an idle map.

\begin{corollary}[Effective finite obstructions]\label{cor:effective}
In dimension two, for explicitly given rational or real-algebraic
polygons $S,D$ and volume-preserving linear maps, every nonempty
$\delta_M$ is a computable real-algebraic number. If the hypothesis of
Theorem~\ref{thm:volume} holds, one can compute a rational
$0<c_M\le\delta_M$ and hence a finite horizon excluding any specified
uniform width $K$ with $M=2^K$.
\end{corollary}
\begin{proof}
Use $M$ pairs of vertex coordinates in $D$ as variables. Containment of
each vertex of $S$ in their hull is expressed with nonnegative
barycentric variables. A finite partition by determinant signs and
cyclic orders identifies each possible convex hull, including
repetitions and collinear points. On every such piece the areas of
$P$ and $\conv\bigcup_LLP$ are polynomial expressions, by the polygon
area formula. Thus the graph of the objective, its feasible domain and
the statement that a number is its minimum are first-order formulas over
the real closed field of algebraic numbers. Quantifier elimination,
followed by real-root isolation, computes the minimum and a positive
rational lower bound when it is positive \cite{BPR}.
Algebraic inputs are represented by defining polynomials and isolating
intervals. The procedure is finite; no polynomial complexity in $K$,
the horizon, or a succinct circuit description is asserted.
\end{proof}
